\subsection{What the evidence supports}
(i) \textbf{As an annual predictor}, the learner adds at most a small, non-robust amount to persistence (\AnnSeedMeanGain\,pp on average over learner seeds, against a persistence error of \ResTVDPersist\,pp), and it is significantly worse than persistence on the wider panel; since it also uses same-year covariates (a nowcast), the honest reading is that last year's allocation is very hard to beat. (ii) \textbf{As a peer-informed prior}, it helps without history on the wide ten-division panel (about 1.3\,pp of a 15\,pp error, surviving Holm correction for the stack and the nearest-economies average), only weakly on the 16-function panel, and for health and R\&D intensity in the WDI panel, where a plain ridge regression already recovers most of the gain; it does not help for defence. (iii) \textbf{As an explanation device}, SHAP and LIME agree on direction more than ranking and both explain an association learned from peers, with ambiguity when macro indicators are correlated (synthetic test). (iv) \textbf{Under shocks}, the model reproduces neither the aggregate direction nor the magnitude of reallocation, so analyst judgment dominates. A system built on such a model should expose this competence map to its users. As a rule of thumb, the output is a \emph{reference point and prompt for questions} when an economy has no history, and at best a sanity check when last year's allocation exists.

\subsection{Intended use and human oversight}
The use cases (gap-filling, peer benchmarking, teaching) are descriptive. A cautious workflow would show last year's allocation first, the model reference as a peer benchmark, the gap in percentage points, and the nearest economies and attributions behind it; the analyst applies statutory bounds and records a rationale, and overrides are logged. The model must never auto-approve, and large gaps should prompt investigation, not correction: a model trained on past allocations reproduces past priorities and past biases. None of this has been tested with real users (Section~\ref{sec:open}).

\subsection{Threats to validity and limitations}
\textbf{Construct validity.} Targets are \emph{executed} general-government expenditure by COFOG function, not budget proposals; COFOG definitions are harmonised but practice differs between countries, the 16 ``functions'' mix COFOG levels (Section~\ref{sec:data}), and WDI sector series follow different definitions (e.g.\ UNESCO education, WHO health accounts) from COFOG. TVD weights large functions more; no log-ratio error is reported. \textbf{Information set.} The annual regime is a nowcast with realised same-year covariates; a strict forecast with lagged covariates, and publication lags of COFOG, were not modelled. \textbf{External validity.} The 16-function panel is OECD/EU-centred (high-income); the wide panel adds 35 middle- and low-income economies via the IMF, but only three are low-income, so low-income conclusions are exploratory. \textbf{Internal validity.} Hyper-parameters were fixed (or tuned within training economies) and a feature-specification choice was made on a validation window; nevertheless several comparisons are within sampling error, the temporal evaluation is a single origin dominated by COVID years, a best-of-several selection effect cannot be excluded for the annual hybrid, and with 36 clusters the intervals are wide and approximate. \textbf{Data handling.} Features (not targets) are imputed in three flagged steps \emph{before} cross-validation: interpolation can use later values and peer medians pool held-out economies, so imputation is transductive; the effect is likely small but was not measured, and fiscal features such as government expense (45.7\% imputed) and debt (69.3\%) are heavily affected (Table~\ref{tab:features}). The World Bank's \emph{current} income classification is applied to all years and is partly an outcome of the same economic history; it defines the peer-mean baseline and the group error tables. Some features (central-government expense, tax revenue, debt) are components of, or financed by, the spending envelope in Panel B and the total-budget module. \textbf{Baselines.} The peer mean has four classes; region-by-income and GDP-per-capita baselines were not run. \textbf{Explanations.} SHAP/LIME values depend on the background data and on correlated features; LIME is stochastic and its local fit is modest; the deletion test is weak; the explanations in the dashboard are for the XGBoost member, not the blend. \textbf{Deployment artefacts.} For an economy that is part of the training panel, the demonstration API returns in-sample fits (only the nearest-economies list excludes the queried economy), so its accuracy for such economies is higher than the leave-country-out errors reported here. \textbf{Non-stationarity.} Reallocation in crises (COVID-19, 2022 energy shock) is only weakly predicted.

\subsection{Governance, ethics and responsible use}
The tool is for \emph{assisting} analysts and for \emph{teaching} what peers allocate; it must not be used to justify cuts or to bypass legislative processes. Recommendations inherit the biases of historical budgets and of which economies publish COFOG data; errors are two to three times larger for lower-income economies, an equity concern that is documented but not resolved here~\cite{mehrabi2021,obermeyer2019}. No personal data are used. A model card and data sheet accompany the release; we recommend an independent audit before any operational use and a human-subject study of analyst use (future work).

\subsection{Reproducibility and data availability}
The processed panels, result files, scripts (\texttt{pipeline/01}--\texttt{27}) and \LaTeX{} sources are provided in the repository. The result macros, tables and figures of the manuscript are written by \texttt{pipeline/06b}, so those numbers cannot drift from \texttt{results/}; a few numbers in the running text are typed by hand from the same result files (fold-seed-0 values quoted in Section~\ref{sec:ext}, the legacy-audit statistics), and the internal audit check cited in Section~\ref{sec:total} is not part of the released pipeline. The dataset, code and manuscript sources are in the repository named in the Code availability statement.

\subsection{Which economies to include}
Holding out a whole region degrades or reverses performance for some functions (defence), cold-start error is two to three times larger for lower-middle- and low-income economies, which are under-represented, and the applicability score tracks this error. We therefore recommend training on all economies with official functional data, adding low-income economies as data become available, and reporting the applicability score so that users know when a reference is an extrapolation.

\subsection{Open items: human-subject evaluation}\label{sec:open}
All claims about supporting analysts remain untested with real users. A protocol, materials and analysis script for a three-condition between-subjects study (reference only; plus SHAP; plus SHAP, LIME and the applicability warning) are in the repository; running it requires ethics approval and about 160 participants for the primary comparison, which we leave to future work.
